\section*{Reproducibility}
\label{sec:reproducibility}
\addcontentsline{toc}{section}{Reproducibility}

Every number in this paper is produced by a committed script, and no figure or
table was edited by hand. The scripts write the files the paper \verb|\input|s, so
a stale figure cannot survive a regeneration unnoticed.

\begin{center}
\small
\begin{tabular}{@{}llll@{}}
\toprule
Script (\texttt{quadprog\_note/}) & Figure & Table & Used in \\
\midrule
\verb|experiment_qp_identities| & --- & --- & below \\
\verb|experiment_qp_pdas| & \verb|quadprog_pdas|
  & \verb|_pdas_defs| & \S\ref{ssec:pdas} \\
\verb|experiment_qp_sweep| & ---
  & \verb|_sweep_defs| & \S\ref{sec:sweep} \\
\verb|experiment_qp_compare| & \verb|quadprog_compare|
  & \verb|quadprog_compare| & \S\ref{ssec:compare} \\
\verb|experiment_qp_portfolio| & \verb|quadprog_portfolio|
  & \verb|quadprog_portfolio| & \S\ref{sec:portfolio} \\
\bottomrule
\end{tabular}
\end{center}

\noindent
Each row's table entry stands for both the tabular fragment and the
\verb|_defs| file of macros beside it, where the script writes both. The identity
script writes no figure, and the sweep script writes one the paper does not
include: its cost curves are the natural way to inspect the two regimes
\S\ref{sec:sweep} describes in words, so the figure is left in place for a reader
who wants to see them.

\paragraph{The identities are checked against the code.} A proof establishes the
identities of Sections~\ref{sec:factorisation} to~\ref{sec:sweep} in exact
arithmetic; it does not establish that the implementation satisfies them, which
matters because it reaches them through a packed triangular factor addressed by
hand and through library calls writing in place into a view of $J$'s own buffer.
\verb|experiment_qp_identities| evaluates both sides of all \qpIdentCount{} of them
against the solver's own internal state and not a reimplementation:
\qpIdentChecks{} checks over \qpIdentStates{} factorisation states, worst relative
residual $\qpIdentWorst$. Two come out exactly zero, and one of those is
Remark~\ref{rem:scaleexact}.

\paragraph{Software under study.} \texttt{cvx-quadprog}, version
$0.4.2$~\cite{cvxquadprog}, archived with a DOI. The measurements are tied to that
version rather than to the latest release, so that the reference keeps pointing at
the code that produced them. The experiments import the package's private modules
where the claims are about private state, namely $J$, the packed $R$, and the
vectors $d$, $z$, $r$, and they instrument the shipped routines by wrapping them
instead of reimplementing them, so what is measured is what runs.

\paragraph{Regenerating everything.} From the paper's source tree, \verb|make figures|
runs all five scripts and rewrites every generated file; \verb|make compile| then
rebuilds this document. Each script also runs standalone as a module. Setting
\verb|EXPERIMENT_SMOKE=1| shrinks every sweep so that the whole set completes in
seconds, which is how the continuous-integration suite exercises each code path
without reproducing the full runs, and \verb|EXPERIMENT_OUT| redirects the outputs
so that a reduced run cannot overwrite the committed ones.

\paragraph{Determinism.} Every random problem is drawn from a seed derived
arithmetically from the family, size and instance index, so the tables reproduce
across runs and machines. Timings do not, and are not claimed to: they were taken
on one machine against NumPy built on Apple Accelerate, with no thread count set.
That is not incidental. This solver pushes its work into BLAS calls, so its
timings inherit whatever threading decision the installed library makes by default,
and those defaults differ between libraries by more than the effects measured here.
Accelerate exposes no thread control, so there was nothing to set; on a library
that does, the same experiment can read differently. The residuals of
the identity check below and the counts of \S\ref{ssec:pdas} are machine-independent
to the last digit or two of rounding, and are what the paper rests on.

\paragraph{Optional dependency.} The comparison in \S\ref{ssec:compare} includes
the reference C implementation where a wheel for it is installed; building it needs
the C toolchain the package under study exists to avoid, so that row is omitted,
with a printed note, on a host that lacks it. Every other row is unconditional.
